\section{The exact cost of obstructing simultaneous torsion specialization}
\label{sec:torsion-minimax}

We answer Question (ii) of the quantitative-twists report by working in one variable. A uniform polynomial bound in all dimensions would have to hold here. Logarithms are natural unless a base is specified.

Fix a prime $p$, a radial lower index $h_0\ge1$ and $\rho\ge
1/\varphi(p^{h_0})$. Put $N_j=(p-1)p^{j-1}$, $r_j=1/N_j$, and use
$v(p)=1$, $v(0)=+\infty$. The $h$-th shell consists of the points
$z=\zeta-1$ for $\zeta$ of exact order $p^h$. Consider a finite packet of
nonzero polynomials $g_i\in\mathcal O_{\mathbf C_p}[T]$ with witnesses
$x_i$ satisfying $v(x_i)\ge\rho$ and $t_i=v(g_i(x_i))<\infty$. Its
budget is $B=\sum_i t_i$. Its loss at $z$ is
\[
             D(z)=\sum_i\max\{0,v(g_i(z))-t_i\}.
\]
A shell is blocked at tolerance $\delta\ge0$ if every point on it has
$D(z)>\delta$. Define the extended-real function
\[
 S_h(L)=\frac1{N_h}\sum_{j=1}^{h-1}N_j\left\lfloor\frac{L}{N_j}\right\rfloor
 \quad(0\le L<N_h),\qquad S_h(N_h)=+\infty .
\]
Let $m_h$ be the least integer $L\in\{0,\ldots,N_h\}$ for which $S_h(L)>\delta$.
In particular the empty sum at $h=1$ is zero. The strict threshold is
essential when $\delta$ is an attained loss.

\begin{theorem}[Exact obstruction budget]\label{thm:exact-obstruction}
The minimum budget of a finite packet blocking every shell from $h_0$
through $H$ is
\[
              C_{p,\delta}(h_0,H)=\sum_{h=h_0}^{H}\frac{m_h}{N_h}.
\]
This minimum is attained using only monic linear polynomials with a common
witness $0$ and coefficients in the ring of integers of
$\mathbf Q_p(\zeta_{p^H})$. It also holds if $t_i$ is replaced throughout by
the Gauss valuation $G_\rho(g_i)=\min_k(v([T^k]g_i)+k\rho)$.
\end{theorem}

\begin{proof}
For $g=A\prod_\alpha(T-\alpha)$, count roots with multiplicity. Multiplicativity
of the Gauss valuation gives $G_s(g)=v(A)+\sum_\alpha\min(s,v(\alpha))$.
This follows by multiplying the reductions of scaled polynomials at rational
weights, or by choosing the smallest exponent attaining each coefficient
minimum. Integral coefficients give $G_0(g)\ge0$. Therefore
\[
 \sum_{h=h_0}^H r_h\#\{\alpha:v(\alpha)=r_h\}
 \le G_\rho(g)-G_0(g)\le G_\rho(g)\le v(g(x_i))
\]
when $g=g_i$. On the $h$-th shell every root with valuation different
from $r_h$ contributes exactly its Gauss term, by the ultrametric inequality.
Thus for every $z$ on this shell,
\[
 D(z)\le\sum_{\substack{\alpha\text{ across the packet}\\v(\alpha)=r_h}}
                       (v(z-\alpha)-r_h).
\]
Indeed $G_{r_h}(g_i)\le G_\rho(g_i)\le t_i$ and every summand on the
right is nonnegative. This argument works with either stated baseline.

Fix a primitive $p^h$-th root $\xi$. For each of the $L$ roots $\alpha$ on
this shell mark an exponent $b$ such that $\xi^b-1$ is closest to $\alpha$.
Here $b$ is a unit modulo $p^h$, and ties may be resolved arbitrarily. If
$a\ne b$ modulo $p^h$, the nearest-point choice implies
$v(\xi^a-1-\alpha)\le v(\xi^a-\xi^b)$. Writing $k=v_p(a-b)$ gives
\[
 v(\xi^a-\xi^b)-r_h
 =\frac{p^k-1}{N_h}=\frac1{N_h}\sum_{j=1}^k N_j.
\]
The first equality uses $v(\xi^{a-b}-1)=1/N_{h-k}$, from the cyclotomic
polynomial at $1$. Marks are allowed to coincide.

If $L<N_h$, begin with a least occupied unit residue class modulo $p$ and
at every stage choose its least occupied lift modulo the next power of $p$.
There are $p-1$ choices at the first stage and $p$ at every subsequent one.
The resulting class modulo $p^j$ contains at most $\lfloor L/N_j\rfloor$
marks, by induction. At level $h$ it is empty. Summing the displayed distance
identity over marks at this unmarked endpoint $a$ gives
$D(\xi^a-1)\le S_h(L)$. Consequently every blocked shell has at least
$m_h$ roots. The Gauss-budget inequality, summed over all polynomials,
establishes the lower bound in the theorem.

For equality, on level $h$ mark the first $m_h$ positive integers prime to
$p$ in their natural order. For each marked $b$, insert the linear
polynomial $T-(\xi^b-1)$ in the packet. Its valuation at $0$ and its
$\rho$-Gauss valuation are both $r_h$. In successive blocks of $N_j$ unit
exponents, every unit residue class modulo $p^j$ appears exactly once.
Hence every residue class at level $j$ contains at least
$\lfloor m_h/N_j\rfloor$ marks. A marked endpoint gives an infinite loss.
At any unmarked endpoint the distance identity gives loss at least
$S_h(m_h)>\delta$. This also proves the assertion when $m_h=N_h$ since all
endpoints are then marked. Linear factors from other shells have zero
positive excess there. Taking all these polynomials gives exactly the
claimed cost.
\end{proof}

At each shell one may multiply its linear factors into one polynomial,
since on that shell each individual excess is nonnegative and on every
other shell their excesses have the same sign. Thus at most
$H-h_0+1$ polynomials suffice if degree is unrestricted.
The balanced residue counts are an instance of the fixed-divisor phenomenon
studied in Bhargava's $P$-orderings (M.~Bhargava, \emph{J. Reine Angew.
Math.} 490 (1997), 101--127, \href{https://doi.org/10.1515/crll.1997.490.101}{doi:10.1515/crll.1997.490.101}).
Here the positive-part loss and distribution of the witness budget across
different cyclotomic shells determine the global optimum.

\begin{corollary}[Exponential depth is necessary]\label{cor:exponential-depth}
For fixed $p,h_0$ and $\delta>0$,
\[
 C_{p,\delta}(h_0,H)=\delta\log H+O_{p,h_0,\delta}(1).
\]
The least universal upper shell index given budget $B$ is the first $H$
with $C_{p,\delta}(h_0,H)>B$. Consequently its logarithm equals
$B/\delta+O_{p,h_0,\delta}(1)$ as $B\to\infty$.
\end{corollary}
\begin{proof}
The upper estimate $S_h(L)\le(h-1)L/N_h$ holds for $L<N_h$. For large
$h$ it follows that $m_h/N_h\ge\delta/(h-1)$, also when $m_h=N_h$.
For the reverse bound let $q=h-\lceil 3\log_p h\rceil$; eventually
$1\le q<h$. For any $L<N_h$,
\[
 S_h(L)\ge\frac{qL-\sum_{j=1}^q N_j}{N_h}
          =\frac{qL-p^q+1}{N_h}.
\]
Take $L=\lfloor(\delta N_h+p^q)/q\rfloor+1$, eventually smaller than
$N_h$. This makes $S_h(L)>\delta$. Since $p^q/N_h=O_p(h^{-3})$ and
$h-q=O_p(\log h)$, one gets
$m_h/N_h\le\delta/h+O_{p,\delta}(\log(h)/h^2)$. Summation and convergence
of the error series give the formula; the strict inequality in the exact
criterion gives the asserted least index.
\end{proof}

\paragraph{Zero tolerance.}
For $h\ge2$, $S_h(L)$ first becomes positive at $L=p-1$; the same number
marks the whole first shell. Thus the exact finite obstruction cost at
$\delta=0$ is $\sum_{h=h_0}^{H}p^{1-h}$. For example, at $p=2$ and $h_0=1$,
arbitrarily many shells can be blocked at zero tolerance using budget less
than $2$. This shows why the positive-tolerance hypothesis in the asymptotic
law matters.

